\chapter{模10 BCD计数器}

\section{BCD计数器与二进制计数器的本质区别}

到目前为止，我们讨论的计数器都是\textbf{二进制计数器}：

\begin{itemize}
  \item 数是用\textbf{整个二进制数}表示的
  \item 模12计数器：4位二进制（0000-1011），一个整体
  \item 模60计数器：6位二进制（000000-111011），一个整体
\end{itemize}

\textbf{BCD计数器}完全不同：

\begin{keypoint}[BCD计数器的本质]
BCD（Binary-Coded Decimal）计数器将\textbf{每个十进制位}单独用一个4位二进制计数器来表示。

\begin{itemize}
  \item 每个十进制位 = 一个独立的模10计数器（0-9，4位）
  \item 个位计到9后，产生进位给十位
  \item 十位每收到一个进位，加1
\end{itemize}
\end{keypoint}

\section{一位BCD计数器 = 模10计数器}

一位十进制数字（0-9）需要一个模10计数器。

\begin{enumerate}
  \item \textbf{模值N}：$N = 10$
  \item \textbf{寄存器位宽}：$k = 4$（$\lceil \log_2 10 \rceil = 4$）
  \item \textbf{状态序列}：0000→0001→0010→...→1001(9)→0000(0)
  \item \textbf{清零条件}：cnt = 9(1001)
\end{enumerate}

\section{逐拍状态分析}

\begin{beatanalysis}[模10 BCD计数器]
\begin{center}
\begin{tabular}{|c|c|c|c|}
\hline
\textbf{时钟拍} & \textbf{上升沿后Q} & \textbf{十进制} & \textbf{进位} \\
\hline
1 & 0001 & 1 & 0 \\
2 & 0010 & 2 & 0 \\
3 & 0011 & 3 & 0 \\
4 & 0100 & 4 & 0 \\
5 & 0101 & 5 & 0 \\
6 & 0110 & 6 & 0 \\
7 & 0111 & 7 & 0 \\
8 & 1000 & 8 & 0 \\
9 & 1001 & 9 & 0 \\
\hline
\rowcolor{yellow!20}
\textbf{10} & \textbf{0000} & \textbf{0} & \textbf{1} \\
\hline
11 & 0001 & 1 & 0 \\
\hline
\end{tabular}
\end{center}
\end{beatanalysis}

\textbf{关键观察}：模10计数器在状态9(1001)之后强制归零——而1001+1本应是1010(10)。所以\textbf{状态1010到1111(10-15)}永远不会出现。

\section{VHDL实现}

\begin{lstlisting}[caption={模10 BCD计数器（个位）}]
library ieee;
use ieee.std_logic_1164.all;
use ieee.std_logic_unsigned.all;

entity bcd_counter is
  port (
    clk   : in  std_logic;
    rst_n : in  std_logic;
    count : out std_logic_vector(3 downto 0);
    carry : out std_logic  -- 进位到十位
  );
end bcd_counter;

architecture behavioral of bcd_counter is
  signal cnt : std_logic_vector(3 downto 0);
begin
  process(clk, rst_n)
  begin
    if rst_n = '0' then
      cnt <= (others => '0');
    elsif rising_edge(clk) then
      if cnt = "1001" then    -- = 9
        cnt <= (others => '0');
      else
        cnt <= cnt + 1;
      end if;
    end if;
  end process;

  count <= cnt;
  carry <= '1' when cnt = "1001" else '0';
end behavioral;
\end{lstlisting}

\section{BCD模10 vs 二进制模10}

\begin{center}
\begin{tabular}{|l|c|c|}
\hline
\textbf{特性} & \textbf{二进制模10} & \textbf{BCD模10} \\
\hline
位宽 & 4位 & 4位 \\
状态数 & 10个 & 10个 \\
表示方式 & 二进制（整体） & 一个十进制位 \\
清零条件 & cnt=1001(9) & cnt=1001(9) \\
本质 & \multicolumn{2}{c|}{\textbf{完全相同——都是模10计数器！}} \\
\hline
\end{tabular}
\end{center}

\section{BCD模10 vs 多位BCD计数器}

单个BCD计数器 = 模10计数器。

\textbf{真正的BCD计数器}是指多位BCD级联：个位+十位+百位+...

每个十进制位都是独立的模10计数器，个位向十位进位，十位向百位进位。

\textbf{这将在下一章详细展开。}
